\chapter{Vérification formelle d'identités algébriques finies}
\label{app:lean}

Les égalités finies suivantes sont formalisées dans Lean et décidées par
réduction du noyau logique, sans axiomes externes supplémentaires. Les théorèmes
démontrent :
\begin{longtable}{p{.40\textwidth}p{.50\textwidth}}
\toprule
Théorème & Contenu\\
\midrule
\endhead
\texttt{aw\_symmetric}
& symétrie de la matrice de Paley--Witt \(A_W\) ;\\
\texttt{aw\_square\_minus\_identity}
& \(A_W^2=-I_6\) sur \(\mathbb F_3\) ;\\
\texttt{x\_sq\_add\_one\_has\_no\_root}
& irréductibilité de \(x^2+1\) sur \(\mathbb F_3\) ;\\
\texttt{dual\_has\_nonzero\_nilpotent}
& existence d'un nilpotent dans l'algèbre parabolique ;\\
\texttt{split\_has\_nontrivial\_idempotent}
& existence d'un idempotent non trivial dans l'algèbre hyperbolique ;\\
\texttt{canonical\_charge}
& \(\sum_{m=1}^{12}K_m=6263\);\\
\texttt{hadamard\_orbit\_one/two/three}
& les trois orbites de Hadamard produisent l'état dodécaphasique publié.\\
\bottomrule
\end{longtable}

La formalisation reçoit les matrices et le vecteur dodécaphasique comme données
typées et vérifie leurs conséquences algébriques. La construction de l'état
signé et l'énumération de (104\,976) germes appartiennent aux théorèmes
combinatoires antérieurs et ne sont pas remplacées par cette réduction formelle.
